\begin{proofbox}
\textbf{1. 焦点距离积的解析证明}:
\begin{enumerate}
    \item 设切线方程为 $l: kx - y + m = 0$, 其中满足 $m^2 = a^2k^2 + b^2$.
    \item 焦点坐标为 $F_1(-c, 0)$ 和 $F_2(c, 0)$.
    \item $d_1 = \frac{|-kc + m|}{\sqrt{k^2+1}}$, $d_2 = \frac{|kc + m|}{\sqrt{k^2+1}}$.
    \item $d_1 \cdot d_2 = \frac{|m^2 - k^2c^2|}{k^2+1} = \frac{|a^2k^2 + b^2 - k^2(a^2-b^2)|}{k^2+1} = \frac{b^2k^2 + b^2}{k^2+1} = b^2$.
\end{enumerate}

\textbf{2. 垂足位置性质 (主圆轨迹) 的证明}:
\begin{enumerate}
    \item 设切线 $l: y = kx + m$, 垂足为 $H(x, y)$.
    \item 由于 $H$ 在切线上, 满足方程 (1): $y - kx = m$. 
    \item 过焦点 $F_2(c, 0)$ 作切线的垂线, 其斜率为 $-1/k$, 方程为 (2): $y - 0 = -\frac{1}{k}(x - c)$, 整理得 $ky + x = c$.
    \item 将方程 (1) 和 (2) 分别平方后相加:
    $(y - kx)^2 + (ky + x)^2 = m^2 + c^2$.
    \item 左边展开化简: $y^2 - 2kxy + k^2x^2 + k^2y^2 + 2kxy + x^2 = (1 + k^2)(x^2 + y^2)$.
    \item 右边代入切线判定式 $m^2 = a^2k^2 + b^2$:
    $m^2 + c^2 = a^2k^2 + b^2 + (a^2 - b^2) = a^2k^2 + a^2 = a^2(1 + k^2)$.
    \item 故 $(1 + k^2)(x^2 + y^2) = a^2(1 + k^2)$, 即 $x^2 + y^2 = a^2$.
    \item 结论证毕: 垂足 $H$ 的轨迹是以原点为圆心, 长半轴 $a$ 为半径的圆.
\end{enumerate}
\end{proofbox}